\section{Related Work}

\subsection{Benchmark Reliability}

Recent work shows that benchmarks differ substantially in how reliably they guide language model development. \citet{heineman_signal_2025} introduce a signal-to-noise ratio (SNR) framework where signal measures a benchmark's ability to separate models and noise captures score variability across checkpoints or stochastic runs. They show that high-SNR benchmarks better preserve rankings and reduce scaling-law prediction error. Complementary work studies benchmark quality through ranking consistency, discriminability, and capability-alignment deviations \citep{qian_benchmark2_2026}, while IRT-based methods show that many benchmark items have weak discriminative power and can often be replaced by smaller, more informative subsets \citep{zhou_lost_2026,li_adaptive_2025}.

Efficiency-oriented work reaches similar conclusions.
\citet{perlitz_efficient_2024} quantify the reliability cost of benchmark design decisions and show that large suites can be cut by two orders of magnitude with little loss in ranking reliability. \citet{polo_tinybenchmarks_2024} use item response theory to select small representative subsets, estimating MMLU performance from 100 curated examples.
\citet{gupta_improving_2025} propose SMART filtering to remove easy, contaminated, or redundant examples, reducing evaluation size while preserving agreement with independent model rankings. Scaling-law studies further suggest that only some evaluations provide predictable development signal: \citet{chen_scaling_2025} forecast downstream performance from smaller models, while \citet{schaeffer_pretraining_2026} study scaling laws for generative metrics such as pass@$k$. Together, these papers motivate our focus on identifying which multilingual measurements remain reliable across model sizes and nearby pretraining checkpoints, rather than treating all benchmark averages as equally meaningful.

\subsection{Multilingual Evaluation}

Multilingual benchmarks inherit reliability problems from their English counterparts and add new ones. Most popular benchmarks are machine-translated from English, which introduces translation artifacts, cultural mismatch, and label noise that distort cross-lingual rankings \citep{plaza_spanish_2024, singh_global_2024, romanou_include_2024}. Tokenizer performance varies by an order of magnitude across languages, inflating both training and evaluation cost for non-Latin scripts and morphologically rich languages \citep{petrov_language_2023, ahia_all_2023}. Cross-lingual transfer further couples decisions across languages, since adding data for one language can help or hurt others depending on model capacity and language similarity \citep{chang_when_2023, longpre_atlas_2026}. These factors suggest that English reliability estimates do not transfer directly across languages, and motivate a multilingual-specific framework built on controlled pretraining runs.
